\section{Scope, theorem inputs, and the central reductions}

The main unrestricted tasks remain open in this manuscript: an explicit polynomial-size hitting set for general width-$w$ ROABPs, a uniformly efficient positive rule for all rectangular and plethysm multiplicities, and a determinant-versus-padded-permanent separation in the required asymptotic regime. Counterexamples below refute specified transfers, not those open problems. Every stated deduction has an identifier. The accompanying event register distinguishes complete proof readings, finite exact computations, dependency checks, and carried background. No Lean reconstruction or formal verification of the new mathematics is claimed.

The principal positive results are exact rather than heuristic. R-01--R-09 identify the CUE character and ratios operations, positive projector formulas, orbit-closure Gram ranks, and the missing quotient numerator. R-10 closes the height-two rectangle and quadratic determinant case. G-05 excludes every at-most-three-row multiplicity obstruction. G-07/R-12 give an exact four-row permanent-minus-determinant shifted-plethysm formula at size four. R-13 classifies quaternary permanent dominance by exact certificates for sizes one through sixteen and a dimension exclusion thereafter. R-15--R-22 give algebraic dimension growth and an unbounded range of genuine padding. R-22 proves $4m^2-20m^{3/2}+2\leq\dim\mathcal P_{m,4}\leq4m^2-2m+2$ for $m\geq99$, reaching every fixed padding ratio below $\sqrt2$ with an explicit threshold. G-10 makes the resulting degree threshold explicit. G-11 identifies rectangular outer highest weights and an explicit degree progression for these gaps. R-20 first closes the eventual rectangular range, and R-22 lowers its uniform threshold to $m\geq2025$: every integer $m<n<4m/3$ has such gaps, while G-12 excludes the complementary range. This is the exact eventual padding range for four-row outer rectangles. G-13 bounds the first possible rectangular gap by a degree with logarithm $O(n^2\log n)$. R-22 certifies nonrectangular four-row gaps beyond that boundary throughout every fixed ratio below $\sqrt2$, eventually. G-16 specifies such a weight throughout the dimension-certified range: $\lambda=(d(n-m)+md/4,md/4,md/4,md/4)$, with explicit first-degree and progression bounds. An auxiliary dual line removes the rectangular stability restriction. More generally, G-17 classifies every four-row stretched direction under the same dimension hypothesis: for $\nu\vdash ns$, some stretch has a positive gap exactly when $\nu_4>0$ and $\nu_1>(n-m)s$. It gives effective first-stretch and progression bounds through a partial-flag construction. The first successful stretch has polynomial binary length in $n$ and $\log(s+1)$, and logarithm $O_s(n^2\log n)$ for fixed $s$. G-18 strengthens this to every sufficiently large integer stretch, including rectangular directions. It gives an all-integer multiplicity ratio with exponent $\dim\mathcal P_{m,4}-2n^2+1$, an effective converse range, and eventual absence of occurrence obstructions along each fixed four-row direction. G-15 strengthens the necessary first-row condition to $\lambda_1\geq d(n-m)+4$, using a dominant third-jet map, and excludes full fourth-jet dominance by dimension. R-21 gives an exact orbit-multiplicity algorithm for each fixed ambient dimension, with numerical degree rather than bit-length complexity. G-08--G-09 give explicit low-row exclusion bounds at larger padding.

On isolation, I-01 gives the genuine Hoza transfer for independent vertex variables. I-03--I-04 identify coefficient flags and projected matroid cones; I-14 gives the weaker joint-span condition. The rank-profile construction I-21 and endpoint refinement I-30 supply an explicit six-pair interpolation hitting generator for width at most three in one unknown variable order, in every characteristic with the stated field-node convention. I-29 treats all words having one fixed singular interval rank. The full-matrix cone proof I-31 gives polynomial-size hitting sets for every fixed width and a pair budget controlled by the number of distinct interval ranks. Complete proofs are given here; the rank-profile and cone mechanisms are credited to a recent unrefereed preprint, whose broader affine-orbit claims are not assumed. I-32 refutes a polynomial coordinate-cone deletion bound even at a common center with every interval rank preserved. I-33 instead gives explicit small prime-cocharacter families for two dense products, their exact limitation, and a scalar product-closure repair. I-34 extends the prime family to bounded local nullity among separated products, with an exact three-versus-four-term cylinder distinction. P-14 uses scalar bridges to give $w-1$ interpolation pairs for factors of rank one or full rank, including one pair at width two. The unresolved target is a bound polynomial jointly in the growing width, length and degree; even a polynomial pair budget does not by itself meet that sampling target. The precise conditional matrix-invariant consequence remains I-08.

Purity gives sparse constructions, an actual-sheaf Gram criterion with full conductor accounting, and limitations of energy approximation. P-11 identifies explicit common arithmetic evaluation points with hitting-set encodings in its rational model. P-12 constructs such a point for the width-three class, together with an integer hitting curve and an exact trace-dual coefficient-span map. P-13 states the general centered-interpolation encoding and reduces the width-three field degree to $(nD+1)^{13}$ using I-30. It applies directly to the rank-spectrum construction I-31 with degree $(nD+1)^{2K_{\mathcal R}+1}$. These constructions use elementary algebraic number theory and require no additional RH input.

\subsection{Stipulated analytic inputs and supplied material}

We accept Alp\"oge's RH and Dirichlet GRH as theorem inputs, exactly as requested. In \texttt{ALPOGE\_Main.lean}, the exported declarations are \texttt{riemannHypothesis}, \texttt{riemannZeta\_eq\_zero\_re}, \texttt{LFunction\_ne\_zero\_of\_half\_lt\_re}, \texttt{LFunction\_eq\_zero\_re}, and \texttt{LFunction\_eq\_zero\_re\_of\_isPrimitive}; the predicates are defined in \texttt{ALPOGE\_Statements.lean}. The attached note supplies narrative descriptions of its numbered steps. Neither the imported Lean dependency tree nor the claimed all-files scan is in this attachment, and neither was rebuilt in this run. This is an input convention, not a review of the theorem's truth.

Masullo's general number-field Hecke GRH is used only conditionally, as the prompt's attribution to section~1.2. The supplied statement index has no numbered section~1.2 theorem or its proof. Accordingly we give no invented locator. Its Theorem~13.3.1.1, indexed PDF page~522, is a trace-formula claim; Conjecture~14.1.2.1, indexed page~529, concerns automorphic $\mathrm{GL}_n$ GRH and simplicity and remains explicitly conjectural there. These are different statements. The zero-statistics section names the precise hypotheses discharged by the stipulated input and those that remain.

The archive was unpacked once. The source manifest records the actual supplied files and hashes. The five cote~133 files are LaTeX transcriptions, not original page images. Assertions in the accompanying earlier study about visual readings remain credited to that study; our transcription table records mathematical necessity separately from paleographic confidence.

\subsection{Why the three multiplicity problems are distinct}

Let $V=\mathbb C^{n^2}$, $W=\operatorname{Sym}^nV^*$, and $A=\operatorname{Sym}(W^*)$. For a homogeneous form $f\in W$, write
\[
 m_f(\lambda,d)=\operatorname{mult}_{S_\lambda V}\mathbb C[\overline{GL(V)f}]_d,
 \qquad a_\lambda(d[n])=\operatorname{mult}_{S_\lambda V}A_d.
\]
The symbol $g(\lambda,\mu,\nu)$ denotes the ordinary Kronecker coefficient. One must distinguish the ambient plethysm multiplicity, the connected-stabilizer rectangular Kronecker multiplicity, and the coordinate-ring multiplicity of an orbit closure. The transpose component of the determinant stabilizer and the orbit boundary are separate restrictions. R-03 gives finite counterexamples to identifying them. R-05 computes the closure multiplicity by an actual rank, and R-08 expresses the same requirement as a quotient numerator.

If $X\subseteq W$ is $GL(V)$-stable and $V_k$ is spanned by the first $k$ basis vectors, let $X_k$ be the closure of its restriction image in $\operatorname{Sym}^nV_k^*$. For $\ell(\lambda)\leq k$, the multiplicities in $\mathbb C[X]$ and $\mathbb C[X_k]$ agree. The short verification used repeatedly below is as follows. A highest-weight polynomial of weight $(\lambda_1,\ldots,\lambda_k,0,\ldots,0)$ uses only coefficient variables supported on $V_k$, since all generator weights are nonnegative. It vanishes on $X$ exactly when it vanishes on $X_k$. Taking the quotient of the common ambient highest-weight space proves the assertion. This is a coefficient-subalgebra version of the usual inheritance principle, proved here in precisely the form needed for the orbit restrictions.

The exact identities alone impose no efficient complexity bound. In particular, a nonnegative squared-norm formula may have exponentially many terms; a rank formula may require exponentially large matrices; and a positive dimension may become zero after highest-weight projection. The size and sign statements in each result are therefore part of the theorem, not suppressed implementation details.

\subsection{How the specific new analytic tools enter}

Alp\"oge's step~2 (p.~3) concerns theta derivatives and inert-prime divisibility. It does not, in the supplied four-page text, state a bound for a coefficient-matroid minor. Steps~3 and~5 (p.~3) motivate trace and square-root-cancellation arguments; P-02, P-04 and P-06 implement the corresponding curve/sheaf constructions with their exact degree and conductor budgets. Steps~4 and~7 (pp.~3--4) concern determinant and Gan--Gross--Prasad period positivity for specified groups; R-03 identifies why that branching problem cannot be substituted for the tensor subgroup without a proved functorial bridge. Step~8 (p.~4) concerns theta-lift multiplicities; I-03 and I-14 identify the different coefficient-span quantifiers a proposed isolation transfer must preserve.

Step~13 (p.~4) invokes a Bost slope inequality, and step~17 (p.~4) concerns the displayed arithmetic cohomology groups. The supplied note gives no reusable estimates for the ranks, heights or valuations of the ROABP coefficient maps used here. Their possible use therefore has a concrete prerequisite: construct the relevant arithmetic bundle or cohomology object, identify the desired minor or quotient with its invariant, and retain effective size bounds. This is a dependency statement about the supplied inputs, not a claim that such a construction is impossible. The integral and special-fiber analyses below identify where that next construction would act.
